\documentclass[12pt]{article}

\usepackage[utf8]{inputenc}
\usepackage[T1]{fontenc}
\usepackage[spanish]{babel}
\addto\captionsspanish{\renewcommand{\tablename}{Tabla}}
\usepackage{amsmath, amssymb}
\usepackage{geometry}
\usepackage{booktabs}
\usepackage{array}
\usepackage{enumitem}

\geometry{margin=2.5cm}
\renewcommand{\arraystretch}{1.2}
\newcolumntype{L}[1]{>{\raggedright\arraybackslash}p{#1}}
\newcolumntype{C}[1]{>{\centering\arraybackslash}p{#1}}

\begin{document}

\section*{Modelamiento 1: Distribución de Frecuencias y Mediana}

\textit{Paso 1. Contar con la muestra y determinar su tamaño.} Se cuenta cuántas respuestas hay en total. Ejemplo: situación laboral, $n=148$ titulados.

\textit{Paso 2. Contar la frecuencia de cada respuesta ($f_i$) y dividir entre $n$ para obtener la relativa ($f_i/n$).}

\begin{tabular}{lccc}
\toprule
Situación laboral & $f_i$ & $f_i/n$ & \% \\
\midrule
Trabaja en una organización & 108 & $108/148=0.730$ & 73.0 \\
Emprendimiento propio & 15 & $15/148=0.101$ & 10.1 \\
No trabaja & 25 & $25/148=0.169$ & 16.9 \\
\midrule
Total & 148 & & 100.0 \\
\bottomrule
\end{tabular}

\textit{Lectura:} la más frecuente es ``Trabaja en una organización'' (73.0\%). Si se quiere una sola proporción agrupada (ej. tasa de empleo), se suman las frecuencias de las categorías que correspondan y se divide entre $n$:
\[
\text{Tasa de empleo} = \frac{108+15}{148} = \frac{123}{148} = 0.8311 = 83.1\%
\]

\textit{Paso 3. Mediana (cuando la media puede engañar).} Si unos pocos valores muy altos elevan la media, se calcula también la mediana, que no se deja arrastrar por esos valores extremos.

\begin{enumerate}
  \item Ordenar los $n$ datos de menor a mayor.
  \item Ubicar el centro: si $n$ es impar, la posición $\frac{n+1}{2}$; si es par, el promedio de las posiciones $\frac{n}{2}$ y $\frac{n}{2}+1$.
\end{enumerate}

Ejemplo (años de vida profesional, $n=143$, impar): posición central $=\frac{143+1}{2}=72 \rightarrow$ Mediana $=6$ años, frente a una media de $8.95$ años. La media es mayor porque unos pocos titulados con muchos años de experiencia la elevan; la mediana describe mejor al titulado típico. Por eso sirve como punto de corte (ej. junior: hasta 6 años; consolidado: más de 6 años).

\section*{Modelamiento 2: Intervalo de Confianza para una Media}

Sirve para estimar, a partir de una muestra, entre qué dos valores está la media de \textbf{toda} la población (no solo de quienes respondieron), con una confianza del 95\%. Se usa cuando no se conoce la varianza poblacional (Coss Bu, cap. 6).

Ejemplo: años de vida profesional, $n=143$.

\textit{Paso 1. Media muestral.}
\[
\bar{x} = \frac{\sum f_i x_i}{n} = \frac{1280.50}{143} = 8.9545 \text{ años}
\]

\textit{Paso 2. Desviación estándar muestral} (qué tan dispersos están los datos respecto a $\bar{x}$; ecuación 6.5):
\[
S = \sqrt{\frac{\sum f_i (x_i-\bar{x})^2}{n-1}} = \sqrt{\frac{8878.9545}{142}} = 7.9075
\]

\textit{Paso 3. Estadístico teórico de tabla.} Con $\alpha=0.05$ y $n-1=142$ grados de libertad (más de 30, se aproxima con la normal): $t_{142,\,0.025} \approx z_{0.025} = 1.960$.

\textit{Paso 4. Margen de error} (se multiplica el error estándar de la media por el valor de tabla, para saber cuánto puede alejarse $\bar{x}$ del verdadero $\mu$):
\[
\frac{S}{\sqrt{n}}\,t_{n-1,\,\alpha/2} = \frac{7.9075}{\sqrt{143}}\times 1.960 = 0.6613 \times 1.960 = 1.2961
\]

\textit{Paso 5. Construir el intervalo} (restando y sumando el margen de error a la media):
\[
8.9545 - 1.2961 \;\le\; \mu \;\le\; 8.9545 + 1.2961 \;\Longrightarrow\; 7.66 \le \mu \le 10.25
\]

Con 95\% de confianza, la media real (de todos los titulados) está entre 7.66 y 10.25 años.

\section*{Modelamiento 3: Prueba de Independencia Chi-Cuadrada (Tabla de Contingencia)}

Sirve para saber si dos variables categóricas están relacionadas, o si las diferencias que se ven en la tabla se deben al azar de la muestra. Extiende la prueba de las frecuencias de la asignatura ($FE=n/k$, una variable) a dos variables cruzadas.

Ejemplo: situación laboral ($X$) $\times$ satisfacción con la formación ($Y$), $N=147$.

\textit{Paso 1. Hipótesis.} $H_0$: $X$ y $Y$ son independientes. $H_1$: existe dependencia.

\textit{Paso 2. Tabla de contingencia} (frecuencia observada $FO_{ij}$: cuántos respondieron a la vez la fila y la columna):

\begin{tabular}{L{3.6cm}C{2.6cm}C{2.6cm}|c}
\toprule
$X \backslash Y$ & En desacuerdo & De acuerdo & Total \\
\midrule
Trabaja en organización & 27 & 81 & 108 \\
Emprendimiento propio & 5 & 9 & 14 \\
No trabaja & 6 & 19 & 25 \\
\midrule
Total & 38 & 109 & $N=147$ \\
\bottomrule
\end{tabular}

\textit{Paso 3. Frecuencia esperada} (lo que habría en cada celda si $X$ e $Y$ \textbf{no} estuvieran relacionadas -- generalización de $FE=n/k$ para totales de fila/columna no uniformes):
\[
FE_{ij} = \frac{(\text{total fila})(\text{total columna})}{N}, \quad \text{ej. } FE_{11} = \frac{108\times38}{147} = 27.92
\]

\textit{Paso 4. Estadístico} (sumar, en cada celda, el cuadrado de la diferencia entre lo observado y lo esperado, dividido entre lo esperado -- así se estandariza para poder comparar contra un patrón, la distribución $\chi^2$):
\[
X_0^2 = \sum \frac{(FO_{ij}-FE_{ij})^2}{FE_{ij}} = 0.0407+0.7106+0.0447 = 0.7960
\]

\textit{Paso 5. Estadístico teórico.} $gl=(f-1)(c-1)=(3-1)(2-1)=2$. De tabla: $X^2_{2,\,0.05}=5.991$.

\textit{Paso 6. Comparar y concluir.}
\[
0.7960 < 5.991 \;\Longrightarrow\; \text{no se rechaza } H_0
\]
No hay evidencia de que la satisfacción dependa de la situación laboral.

\section*{Modelamiento 4: Análisis Exploratorio Multivariable}

Sirve para ver un patrón entre \textbf{tres} variables a la vez: dos cualitativas ($X$, $Y$) y una cuantitativa ($Z$). No es una prueba (no hay valor de tabla contra el cual comparar) -- es exploración descriptiva, para decidir si después amerita una prueba de independencia aparte. Sigue el ejemplo de la asignatura (apunte del docente).

Ejemplo: $X$=nivel del posgrado cursado, $Y$=fuente de financiamiento, $Z$=años de vida profesional ($n=96$, solo quienes hicieron formación complementaria).

\textit{Paso 1. Elegir $X$, $Y$ cualitativas y $Z$ cuantitativa.}

\textit{Paso 2. Para cada combinación $X_i$-$Y_j$ que aparece, contar cuántos casos hay ($n_{ij}$) y promediar $Z$ entre ellos ($\bar{Z}_{ij}$):}

\begin{tabular}{llcc}
\toprule
$X$ (nivel) & $Y$ (financiamiento) & $n$ & $\bar{Z}$ (años) \\
\midrule
Diplomado & Recursos propios & 61 & 6.85 \\
Maestría & Recursos propios & 22 & 11.95 \\
Maestría & Beca parcial & 3 & 10.00 \\
Especialidad & Recursos propios & 3 & 12.33 \\
Doctorado & Recursos propios & 2 & 16.00 \\
Diplomado & Beca parcial & 2 & 6.50 \\
Maestría & Beca total & 2 & 25.50 \\
Diplomado & Beca total & 1 & 13.00 \\
\midrule
\multicolumn{2}{l}{Total / Promedio general} & 96 & 8.93 \\
\bottomrule
\end{tabular}

\textit{Paso 3. Para cada categoría de $X$: la $Y$ más frecuente, y si su $\bar{Z}$ queda sobre o bajo el promedio general.} Ejemplo con Doctorado ($n=2$): financiamiento más frecuente ``Recursos propios'' (100\%); $\bar{Z}=16.00$ años, por encima del promedio general (8.93 años).

\textit{Patrón encontrado:} a mayor nivel del programa, mayor la antigüedad promedio de quienes lo cursaron (Doctorado y Maestría por encima del promedio; Diplomado por debajo). Recursos propios domina el financiamiento en todos los niveles.

\end{document}
